\documentclass[11pt,spanish]{article}

% =========================================================
% PLANTILLA BASE PARA INFORMES ACADÉMICOS / TÉCNICOS
% =========================================================

% --- Idioma y codificación ---
\usepackage[spanish,es-tabla]{babel}
\usepackage[T1]{fontenc}
\usepackage[utf8]{inputenc}
\usepackage{textcomp}
\DeclareUnicodeCharacter{20A1}{\textcolonmonetary}
\DeclareUnicodeCharacter{00F7}{\(\div\)}
\DeclareUnicodeCharacter{00D7}{\(\times\)}
\DeclareUnicodeCharacter{2212}{\(-\)}

% --- Márgenes / papel ---
\usepackage{geometry}
\geometry{a4paper, margin=2.7cm}

% --- Matemáticas ---
\usepackage{amsmath}

% --- Tablas, enlaces y elementos visuales ---
\usepackage{array}
\usepackage{tabularx}
\usepackage{booktabs}
\usepackage{longtable}
\usepackage{graphicx}
\usepackage{float}
\usepackage{xcolor}
\usepackage{tikz}
\usetikzlibrary{positioning,shapes.geometric,arrows.meta,calc}
\usepackage{enumitem}
\usepackage{ragged2e}
\usepackage[hidelinks]{hyperref}

% --- Encabezado y pie de página ---
\usepackage[myheadings]{fancyhdr}
\usepackage{lastpage}
\setlength{\headheight}{30pt}
\setlength{\footskip}{30pt}

% --- Interlineado ---
\linespread{1.2}
\sloppy

% =========================================================
% DATOS DEL DOCUMENTO
% =========================================================
\newcommand{\Institucion}{Instituto Tecnológico de Costa Rica}
\newcommand{\Campus}{Campus Tecnológico Central Cartago}
\newcommand{\DocumentoTitulo}{Entregable 1 - Business Case}
\newcommand{\DocumentoSubtitulo}{Plataforma Universitaria Integrada:\\[0.2cm] Modernización y Centralización de Procesos Académicos y Administrativos}
\newcommand{\Curso}{Administración de proyectos}
\newcommand{\Profesor}{Alicia Marcela Salazar Hernández}
\newcommand{\FechaDocumento}{6 de junio de 2026}
\newcommand{\PeriodoAcademico}{I Semestre}
\newcommand{\AnioDocumento}{2026}

\newcommand{\AutoresPortada}{%
    Mauricio González Prendas: 2024143009\\
    Isaac Jiménez Blanco: 2024202804\\
    Tamara Robles Camacho: 2024099342%
}

% Datos que aparecen en el encabezado.
\newcommand{\DocHeaderLeft}{}
\newcommand{\DocHeaderRight}{Business Case}
\newcommand{\DocFooterRight}{Página~\thepage\ de~\pageref{LastPage}}

% Metadatos PDF.
\title{\DocumentoTitulo}
\author{\AutoresPortada}
\date{\FechaDocumento}

% =========================================================
% CONFIGURACIÓN DE TABLAS Y UTILIDADES
% =========================================================
\newcolumntype{Y}{>{\RaggedRight\arraybackslash}X}
\renewcommand{\arraystretch}{1.22}
\setlength{\LTpre}{0.5em}
\setlength{\LTpost}{0.5em}

% Imagen segura: si el archivo no existe, muestra un recuadro de marcador.
\newcommand{\SafeImage}[2][0.92\textwidth]{%
    \IfFileExists{#2}{%
        \includegraphics[width=#1]{#2}%
    }{%
        \fbox{%
            \begin{minipage}[c][4cm][c]{#1}
            \centering
            Imagen pendiente:\\
            \texttt{#2}
            \end{minipage}%
        }%
    }%
}

% Figura reutilizable.
\newcommand{\EvidenceFigure}[3]{%
    \begin{figure}[H]
    \centering
    \SafeImage{#1}
    \caption{#2}
    \label{#3}
    \end{figure}%
}

% =========================================================
% ENCABEZADO Y PIE
% =========================================================
\pagestyle{fancy}
\fancyhf{}
\fancyhead[L]{\DocHeaderLeft}
\fancyhead[R]{\DocHeaderRight}
\renewcommand{\headrulewidth}{0.4pt}
\renewcommand{\footrulewidth}{0.4pt}
\fancyfoot[R]{\DocFooterRight}

% Estilo equivalente para páginas horizontales.
\fancypagestyle{widefancy}{%
    \fancyhf{}%
    \fancyhead[L]{\DocHeaderLeft}%
    \fancyhead[R]{\DocHeaderRight}%
    \renewcommand{\headrulewidth}{0.4pt}%
    \renewcommand{\footrulewidth}{0.4pt}%
    \fancyfoot[R]{\DocFooterRight}%
}

% =========================================================
% PÁGINAS HORIZONTALES PARA TABLAS ANCHAS
% =========================================================
\makeatletter
\newcommand{\SyncPageBuilderDimensions}{%
    \global\vsize=\textheight
    \global\@colroom=\textheight
    \global\@colht=\textheight
}
\makeatother

\newenvironment{widepage}{%
    \clearpage
    \hypersetup{pageanchor=false}%
    \paperwidth=297mm
    \paperheight=210mm
    \pdfpagewidth=297mm
    \pdfpageheight=210mm
    \newgeometry{left=2.7cm,right=2.7cm,top=2.7cm,bottom=2.7cm}%
    \setlength{\textwidth}{243mm}%
    \setlength{\textheight}{156mm}%
    \setlength{\linewidth}{\textwidth}%
    \setlength{\columnwidth}{\textwidth}%
    \hsize=\textwidth
    \setlength{\headwidth}{\textwidth}%
    \SyncPageBuilderDimensions
    \pagestyle{widefancy}%
}{%
    \clearpage
    \paperwidth=210mm
    \paperheight=297mm
    \pdfpagewidth=210mm
    \pdfpageheight=297mm
    \restoregeometry
    \setlength{\textwidth}{156mm}%
    \setlength{\textheight}{243mm}%
    \setlength{\linewidth}{\textwidth}%
    \setlength{\columnwidth}{\textwidth}%
    \hsize=\textwidth
    \setlength{\headwidth}{\textwidth}%
    \SyncPageBuilderDimensions
    \hypersetup{pageanchor=true}%
    \pagestyle{fancy}%
}

% =========================================================
% PORTADA
% =========================================================
\newcommand{\MakeCover}{%
\begin{titlepage}
\thispagestyle{empty}
\centering

{\bfseries \huge \Institucion \par}
\vspace{0.25cm}
{\LARGE \Campus \par}

\vfill

{\huge \DocumentoTitulo \par}
\vspace{0.4cm}
{\Large \DocumentoSubtitulo \par}

\vfill

{\bfseries \LARGE Profesora: \par}
{\Large \Profesor \par}

\vfill

{\bfseries \LARGE Estudiantes: \par}
{\Large \AutoresPortada \par}

\vfill

{\bfseries\LARGE Fecha: \par}
{\Large \FechaDocumento \par}

\vfill

{\LARGE \PeriodoAcademico \par}
{\LARGE \AnioDocumento \par}

\end{titlepage}%
}

% =========================================================
% DOCUMENTO
% =========================================================
\begin{document}
\hypersetup{pageanchor=false}

% =========================
% Portada
% =========================
\MakeCover

% =========================
% Índice
% =========================
\pagenumbering{roman}
\pagestyle{empty}
\tableofcontents
\clearpage
\pagenumbering{arabic}
\hypersetup{pageanchor=true}
\pagestyle{fancy}

% =========================
% Contenido
% =========================
\section{Situación actual}

La Universidad Tecnológica La Mejor es una institución superior que actualmente presenta muchos problemas a causa de como gestionan sus procesos académicos, administrativos y financieros, ya que es mediante una serie de sistemas independientes que no se encuentran integrados entre sí. Esta segmentación tecnológica representa un obstáculo para la eficiencia operativa de la institución.

Los sistemas actualmente en uso son: hojas de cálculo para el registro de calificaciones, correo electrónico para la comunicación entre departamentos, sistemas propietarios sin conexión para la gestión financiera y registros físicos o digitales dispersas para el historial académico estudiantil.

\subsection{Diagnostico operativo}

El análisis de la situación actual muestra los siguientes hallazgos críticos:

\begin{itemize}[leftmargin=*]
    \item Duplicidad de datos: Mucha de la información se registra de manera manual en múltiples sistemas, generando inconsistencias y errores.
    \item Falta de visibilidad: Las autoridades no disponen de indicadores consolidados en tiempo real para la toma de decisiones estratégicas.
    \item Procesos lentos: La matrícula estudiantil, el registro de calificaciones y los tramites financieros dependen de intervención manual en cada etapa.
    \item Comunicación deficiente: La información no fluye de forma automática entre el departamento académico, administrativo y financiero.
    \item Insatisfacción estudiantil: los estudiantes deben acudir presencialmente o contactar múltiples oficinas para realizar gestiones básicas.
\end{itemize}

\section{Problema identificado}

El problema central de la Universidad Tecnológica La Mejor se basa en la ausencia de una plataforma integrada que centralice y automatice los procesos académicos, administrativos y financieros de la institución. Esta carencia genera ineficiencias operativas medibles, deterioro en la experiencia del usuario y limitaciones significativas para la toma de decisiones basada en datos.

\subsection{Árbol de problemas}

Los efectos identificados como consecuencia directa del problema central son:

\begin{itemize}[leftmargin=*]
    \item Incremento en los tiempos de atención al estudiante por procesos manuales repetitivos.
    \item Alta tasa de errores en registros académicos y financieros por falta de validación automatizada.
    \item Imposibilidad de generar reportes ejecutivos en tiempo real para la toma de decisiones.
    \item Carga administrativa excesiva sobre el personal docente y administrativo.
    \item Pérdida de competitividad institucional frente a universidades con sistemas modernos.
\end{itemize}

\subsection{Impacto cuantificable}

Se estima que la situación actual genera los siguientes impactos negativos:

\begin{itemize}[leftmargin=*]
    \item Tiempo promedio de atención por tramite estudiantil: 45 a 90 minutos presenciales.
    \item Porcentaje de errores en registros académicos manuales: aproximadamente 8-12\%.
    \item Horas trabajadas de los funcionarios semanales destinadas a tareas administrativas redundantes: más de 120 horas por departamento.
    \item Costo de oportunidad estimado por ineficiencias operativas: significativo para la sostenibilidad institucional.
\end{itemize}

\section{Oportunidad de mejora}

La modernización tecnológica de la Universidad Tecnológica La Mejor representa una oportunidad para transformar su modelo operativo mejorar la experiencia de todos los funcionarios y estudiantes de la institución y posicionarse como una institución de referencia en gestión universitaria en la región.

\subsection{Visión de la solución}

El desarrollo de una Plataforma Universitaria Integrada permitirá consolidar en un único punto de acceso todos los procesos académicos, administrativos y financieros, ofreciendo interfaces diferenciadas para estudiantes, docentes, personal administrativo, área financiera y alta dirección.

\subsection{Oportunidades identificadas}

Automatización de procesos: reducir de forma considerable el tiempo dedicado a tareas manuales y repetitivas:

\begin{itemize}[leftmargin=*]
    \item Centralización de la información: garantizar una única fuente de verdad para todos los datos institucionales.
    \item Mejora en la experiencia estudiantil: brindar un portal accesible para realizar trámites en línea.
    \item Habilitación de inteligencia institucional: generar dashboards ejecutivos con indicadores académicos y financieros en tiempo real.
    \item Reducción de errores operativos: implementar validaciones automáticas que eliminen inconsistencias financieras en tiempo real.
    \item Escalabilidad: diseñar una arquitectura que permita incorporar nuevos módulos o funcionalidades en el futuro.
\end{itemize}

\section{Beneficios esperados}

Para este nuevo sistema, se espera varios beneficios que mejorarán la experiencia de todos los actores de la institución.

\subsection{Beneficios cuantitativos}

En la siguiente tabla se muestra los beneficios junto con la situación actual y la meta esperada:

\begin{table}[H]
\centering
\small
\begin{tabularx}{\textwidth}{p{5.2cm} p{4.1cm} Y}
\toprule
\textbf{Beneficio} & \textbf{Situación Actual} & \textbf{Meta Esperada} \\
\midrule
Tiempo promedio por trámite estudiantil & 45 - 90 min presenciales & 5 - 10 min en línea \\
Tasa de error en registros académicos & 8\% - 12\% & \textless{} 1\% \\
Horas administrativas por semana & \textgreater{} 120 horas/departamento & \textless{} 40 horas/departamento \\
Satisfacción estudiantil (encuesta) & No medida / baja & \textgreater{} 85\% satisfacción \\
Disponibilidad de reportes ejecutivos & 72 hrs promedio & Tiempo real \\
Acceso a trámites fuera del campus & No disponible & 24/7 desde dispositivos inteligentes \\
\bottomrule
\end{tabularx}
\end{table}

\subsection{Beneficios cualitativos}

\begin{itemize}[leftmargin=*]
    \item Mejora en la imagen institucional y percepción de modernidad por parte de estudiantes y docentes.
    \item Mayor transparencia en los procesos académicos y financieros para todos los actores.
    \item Fortalecimiento de la cultura organizacional orientada a la eficiencia y la innovación.
    \item Cumplimiento de estándares de calidad universitaria y posibilidad de acreditación institucional.
    \item Retención estudiantil mejorada gracias a una experiencia universitaria más fluida y satisfactoria.
    \item Reducción del estrés operativo del personal administrativo y docente.
\end{itemize}

\section{Alternativas evaluadas}

Para atender la problemática identificada, se evaluaron tres alternativas estratégicas con distintos niveles de inversión, riesgo y beneficio. La primera alternativa fue a status quo, es decir, dejar el proyecto a como lo están manejando actualmente, la segunda alternativa fue un sistema ERP (Planificación de Recursos Empresariales) Comercial que centraliza la gestión de los procesos operativos, financieros y de recursos humanos de una empresa en una única plataforma y la última alternativa fue un desarrollo a la medida donde se puede adaptar 100\% el software a la medida de la institución. A continuación, se presenta el análisis comparativo:

\begin{widepage}
{\small
\setlength{\tabcolsep}{3pt}
\begin{longtable}{p{4.1cm} p{4.8cm} p{5.8cm} p{7.2cm}}
\toprule
\textbf{Criterio} & \textbf{Alternativa A Status Quo} & \textbf{Alternativa B Sistema ERP Comercial} & \textbf{Alternativa C Desarrollo a la Medida (Recomendada)} \\
\midrule
\endfirsthead
\toprule
\textbf{Criterio} & \textbf{Alternativa A Status Quo} & \textbf{Alternativa B Sistema ERP Comercial} & \textbf{Alternativa C Desarrollo a la Medida (Recomendada)} \\
\midrule
\endhead
Costo de implementación & Ninguno & Muy alto (más de 65 millones) & Moderado (según alcance) \\
Tiempo de implementación & Inmediato & 12 - 24 meses & 4 - 6 meses \\
Adaptación a la institución & Alta (actual) & Baja (genérico) & Muy alta (personalizado) \\
Escalabilidad & Nula & Media & Alta \\
Mantenimiento & Complejo & Dependencia del proveedor & Autonomía interna \\
Mejora operativa esperada & Ninguna & Alta & Muy alta \\
Riesgo de adopción & Bajo & Alto & Medio \\
Recomendación & No viable & No viable (costo) & Opción Seleccionada \\
\bottomrule
\end{longtable}
}
\end{widepage}

\subsection{Justificación de la alternativa seleccionada}

La Alternativa C - Desarrollo a la Medida: fue seleccionada por ofrecer el mejor balance entre costo, tiempo de implementación, adaptabilidad a los procesos específicos de la Universidad y potencial de mejora operativa. A diferencia de un ERP comercial genérico, una solución desarrollada internamente puede ajustarse con precisión a los flujos de trabajo institucionales, evitar licencias costosas y garantizar la autonomía tecnológica a largo plazo.

\section{Justificación económica}

La posibilidad económica del proyecto se sustenta en el análisis de costos de implementación versus los beneficios operativos y estratégicos que generará la plataforma a mediano y largo plazo.

\subsection{Estimación de costos del proyecto}

En la siguiente tabla se muestra la categoría de costo con su respectivo costo estimado y el porcentaje del total del costo final:

\begin{table}[H]
\centering
\small
\begin{tabularx}{\textwidth}{Y p{4.1cm} p{2.7cm}}
\toprule
\textbf{Categoría de Costo} & \textbf{Costo Estimado} & \textbf{\% del Total} \\
\midrule
Mano de obra (equipo de proyecto) & ₡10 millones & 43\% \\
Licencias y herramientas de software & ₡1,700,000 millones & 7\% \\
Infraestructura y hosting & ₡1,900,000 millones & 8\% \\
Capacitación y gestión del cambio & ₡2,300,500 millones & 10\% \\
Pruebas y control de calidad & ₡3,000,000 millones & 13\% \\
Contingencias (10\%) & ₡1,600,000 millones & 8\% \\
Documentación y gestión del proyecto & ₡2,700,000 millones & 11\% \\
TOTAL ESTIMADO & ₡23,200,500 millones & 100\% \\
\bottomrule
\end{tabularx}
\end{table}

\subsection{Retorno sobre la inversión}

Considerando los ahorros operativos estimados derivados de la reducción de horas administrativas, la disminución de errores en registros y la eliminación de procesos manuales redundantes, se proyecta un retorno de la inversión dentro del primer año de operación del sistema.

\begin{table}[H]
\centering
\small
\begin{tabularx}{\textwidth}{Y p{5cm}}
\toprule
\textbf{Fuente de Ahorro / Beneficio} & \textbf{Estimación Anual (USD)} \\
\midrule
Reducción de horas administrativas redundantes & ₡9,000,000 millones \\
Eliminación de errores y correcciones manuales & ₡4,000,000 millones \\
Reducción de costos en materiales físicos e impresión & ₡2,500,000 millones \\
Mejora en retención estudiantil (estimado) & ₡8,000,000 millones \\
Reducción de tiempo de atención al estudiante & ₡3,500,000 millones \\
AHORRO / BENEFICIO ANUAL ESTIMADO & ₡27,000,000 millones \\
\bottomrule
\end{tabularx}
\end{table}

\subsection{Análisis costo-beneficio}

En la siguiente tabla se muestra un análisis de la inversión entre el costo y el beneficio de implementar este proyecto en la institución, beneficio estimado, periodo de recuperación y el retorno de la inversión estimado en 1 y 3 años.

Las fórmulas usadas para sacar dichos resultados son las siguientes:

Período de recuperación: Inversión ÷ Beneficio anual

ROI año 1: (Beneficio año 1 − Inversión) ÷ Inversión × 100

ROI año 3: (Beneficio acumulado 3 años − Inversión) ÷ Inversión × 100

Relación B/C: Beneficio anual ÷ Inversión (año 1), o Beneficio acumulado ÷ Inversión (año 3)

A continuación, los resultados:

\begin{table}[H]
\centering
\small
\begin{tabularx}{\textwidth}{Y p{5.2cm}}
\toprule
Inversión total del proyecto & ₡ 23,200,500 millones \\
Beneficio anual estimado & ₡27,000,000 millones \\
Período de recuperación de la inversión & \textless{}= 12 meses \\
Retorno sobre la Inversión proyectado (año 1) & 16.4\% \\
Retorno sobre la Inversión proyectado (año 3) & \textgreater{}249\% \\
Relación Costo-Beneficio (B/C) & 1.16 (primer año) / \textgreater{}3.50 (año 3) \\
\bottomrule
\end{tabularx}
\end{table}

\section{Conclusión}

La Universidad Tecnológica La Mejor enfrenta un desafío operativo real y urgente: sus procesos académicos, administrativos y financieros operan de forma segmentada, generando ineficiencias medibles, sobrecarga al personal y una experiencia pobre para estudiantes y docentes.

El análisis presentado en este Business Case demuestra con claridad que la inversión en una Plataforma Universitaria Integrada desarrollada a la medida representa la alternativa más favorable, estratégica y rentable para resolver la problemática identificada. Las demás alternativas evaluadas, tanto el mantenimiento del status quo como la adquisición de un ERP comercial, fueron descartadas por su inviabilidad operativa y económica respectivamente.

Los beneficios proyectados, tanto cuantitativos como cualitativos, superan el costo de la inversión inicial, con un período de recuperación estimado menor o igual doce meses y un retorno sobre la inversión superior al 249\% en un horizonte de tres años. Además, la plataforma posicionará a la institución en un nivel de madurez tecnológica que le permitirá competir, crecer y ofrecer servicios de calidad a su comunidad universitaria.

Por lo anterior, se recomienda formalmente la aprobación y ejecución del proyecto Plataforma Universitaria Integrada, autorizando los recursos humanos, tecnológicos y financieros necesarios para su desarrollo en el plazo establecido.

\subsection{Firmas de aprobación}

\begin{table}[H]
\centering
\begin{tabularx}{\textwidth}{Y Y Y}
\toprule
\textbf{Elaborado por} & \textbf{Revisado por} & \textbf{Aprobado por} \\
\midrule
\begin{tabular}[c]{@{}l@{}}Equipo de Proyecto\\Firma y fecha\end{tabular} & \begin{tabular}[c]{@{}l@{}}Project Manager\\Firma y fecha\end{tabular} & \begin{tabular}[c]{@{}l@{}}Rector / Patrocinador\\Firma y fecha\end{tabular} \\
\bottomrule
\end{tabularx}
\end{table}

\end{document}
